\chapter{Анализ предметной области синхронизации игровых состояний}

\section{Общие положения сетевого взаимодействия}

Многопользовательские игровые приложения\cite{smed2002networkgames} представляют собой распределенные системы\cite{tanenbaum2017distributed}, в которых несколько клиентов взаимодействуют с общим игровым миром через сеть. Для обеспечения корректности игрового процесса требуется поддерживать согласованность данных, описывающих состояние мира, действия игроков и результаты игровых событий. Сетевая передача данных сопровождается задержками, потерями пакетов, вариативностью задержки и ограничениями пропускной способности, что и является основными рисками возникновения рассинхронизации между клиентами и сервером.

Наиболее распространённой архитектурой\cite{gamedevnetsync} является модель с выделенным сервером, который выступает в роли центрального компонента, хранящего "истинное" состояние игрового мира. Клиенты направляют серверу информацию о своих действиях и получают обратно обновления состояния.

\section{Игровое состояние и требования к его синхронизации}

Игровое состояние представляет собой совокупность параметров, описывающих положение объектов в игровом мире, их атрибуты и поведение. Чаще всего состояние изменяется дискретно, что упрощает внутреннюю симуляцию.

Процесс синхронизации должен обеспечивать два ключевых свойства:

\begin{itemize}
    \item Согласованность (клиентское представление состояния мира не должно существенно отличаться от серверного).
    \item Отзывчивость (реакция на действия игрока должна быть визуально своевременной, несмотря на сетевые задержки).
\end{itemize}

Взаимное обеспечение этих свойств представляет собой одну из основных проблем сетевых игр, поскольку повышение отзывчивости часто достигается за счёт временных предсказаний, что увеличивает вероятность ошибок и рассинхронизации.

\section{Модели передачи данных и источники рассинхронизации}

Передача данных в многопользовательских играх может осуществляться по нескольким моделям:

\begin{itemize}
    \item Снимки состояния (snapshots) - периодическая отправка сервером полного или частичного состояния объектов.
    \item Событийная репликация (event-based sync) - сервер передаёт только события, влияющие на состояние мира.
    \item Синхронизация по входным данным (input-sync) - клиенты получают входные команды игроков, а симуляция выполняется локально по детерминированным правилам.
\end{itemize}

Каждая модель имеет свои ограничения. Snapshots требуют большого объёма трафика, event-based sync зависит от детерминизма клиентов, а input-sync предъявляет высокие требования к точности вычислений.

Основные источники рассинхронизации включают:

\begin{itemize}
    \item сетевые задержки и её вариативность;
    \item потерю и дублирование пакетов;
    \item недетерминированность клиентской симуляции;
    \item различия в частоте обновления и мощности устройств;
    \item асинхронность обработки событий.
\end{itemize}

\section{Факторы, влияющие на требования к синхронизации}

Требования к точности и частоте синхронизации существенно отличаются в зависимости от жанра игры:

\begin{itemize}
    \item FPS и экшен-игры требуют высокой частоты обновления и минимальной задержки, поскольку ошибки в позиции или времени приводят к критичным изменениям игрового процесса.
    \item MOBA и RPG допускают более низкую частоту синхронизации, поскольку взаимодействия менее чувствительны к маленьким задержкам.
    \item RTS часто используют детерминированную модель мира и обмен только входными данными для масштабируемости.
\end{itemize}

Помимо жанра, важными факторами являются:

\begin{itemize}
    \item количество игроков и объектов в сцене;
    \item нагрузка на сеть и ограничения пропускной способности;
    \item необходимость поддержки слабых клиентских устройств;
    \item плотность взаимодействия игроков друг с другом.
\end{itemize}
